%% toc: P3 24 24.6
\section{Proofs for elementary amalgamation and expansion of the language}
\label{sec:annex-p03-w03-p03-006-amalgamation}

Throughout, \(\Lang(X)\) and the constants \(c_x\) are those of
\cref{def:parameter-expansion}, and \(\operatorname{Diag}\) and
\(\operatorname{ElDiag}\) are the diagrams of \cref{def:diagrams}.
Two facts are used repeatedly.

\emph{Reducts.}  Let \(\Lang\subseteq\Lang^+\), let \(\mathcal C\) be an
\(\Lang^+\)-structure, and let \(s\) assign elements of \(C\) to variables.
For every \(\Lang\)-term \(t\), its value under \(s\) is the same in
\(\mathcal C\) and in \(\mathcal C\mathbin{\upharpoonright}\Lang\), by
induction on \(t\), since the two structures interpret every symbol of
\(\Lang\) identically.  Hence every atomic \(\Lang\)-formula has the same
truth value under \(s\) in both.  The Boolean clauses preserve this, and the
quantifier clauses range over the same domain \(C\).  By induction on
formulas, \(\mathcal C\entails\varphi[s]\) iff
\(\mathcal C\mathbin{\upharpoonright}\Lang\entails\varphi[s]\) for every
\(\Lang\)-formula \(\varphi\).  The same argument shows that a sentence
has the same truth value in any two structures that agree on the symbols
occurring in it and have the same domain.

\emph{Embeddings and quantifier-free formulas.}  Let
\(h:\mathcal M\to\mathcal N\) be an embedding.  The atomic and Boolean cases
of the proof of \cref{lem:iso-invariance} use only that \(h\) is an
embedding, through \cref{lem:term-preservation}, injectivity, and
preservation and reflection of relations.  Hence
\(\mathcal M\entails\chi(\bar a)\) iff \(\mathcal N\entails\chi(h\bar a)\)
for every quantifier-free \(\chi\) and every tuple \(\bar a\) from \(M\).

\emph{Renaming.}  Let \(\mathcal E\) be a structure and
\(k:\mathcal C\to\mathcal E\) an injective map, where \(C\) is a set.  Put
\(Y=E\setminus k(C)\) and define \(\rho\) on \(E\) by \(\rho(k(c))=c\) and
\(\rho(y)=(y,C)\) for \(y\in Y\).  An ordered pair
\((y,C)=\{\{y\},\{y,C\}\}\) is never a member of \(C\): otherwise
\(C\in\{y,C\}\in(y,C)\in C\), a membership cycle excluded by the Axiom of
Foundation.  Distinct \(y\) give distinct pairs, so \(\rho\) is injective.
Transporting the interpretations of \(\mathcal E\) along \(\rho\) gives a
structure \(\rho[\mathcal E]\) with \(\rho\) an isomorphism and
\(C\subseteq\rho[\mathcal E]\).  If \(k\) is an elementary embedding of a
structure \(\mathcal C\), then for every formula \(\varphi\) and tuple
\(\bar c\) from \(C\),
\(\mathcal C\entails\varphi(\bar c)\) iff
\(\mathcal E\entails\varphi(k\bar c)\) iff
\(\rho[\mathcal E]\entails\varphi(\bar c)\), the last step by
\cref{lem:iso-invariance}.  So \(\mathcal C\preccurlyeq\rho[\mathcal E]\),
and for any map \(g\) into \(\mathcal E\), the composite \(\rho\circ g\) is
an embedding, or an elementary embedding, whenever \(g\) is.

\begin{annexitem}{annex:p03-w03-p03-006-expand}{thm:p03-w03-p03-006-expansion-amalgamation}{Amalgamation into an expansion}
\begin{proof}
\emph{Naming.}  Use the constants \(e_c=c_c\) of \(\Lang^+(C)\) for the
elements of \(C\).  For \(b\in B\) define a constant \(\tau(b)\) as
follows.  If \(b=a_i\) for some \(i\in I\), put \(\tau(b)=e_{c_i}\).  This
does not depend on the choice of \(i\): if \(a_i=a_j\), then the atomic
\(\Lang\)-formula \(x_i=x_j\) holds of \(\bar a\) in \(\mathcal B\), so by
hypothesis it holds of \(\bar c\) in
\(\mathcal C\mathbin{\upharpoonright}\Lang\), and \(c_i=c_j\).  If \(b\)
is not a coordinate of \(\bar a\), let \(\tau(b)=d_b\), a new constant; the
\(d_b\) are pairwise distinct and not in \(\Lang^+(C)\).  Let
\(\Lang^*\) be \(\Lang^+(C)\) together with all \(d_b\).  For an
\(\Lang(B)\)-sentence \(\theta\), let \(\theta^\tau\) be the
\(\Lang^*\)-sentence obtained by replacing every constant \(c_b\) by
\(\tau(b)\).  Put
\[
 \Sigma=\operatorname{ElDiag}(\mathcal C)
   \cup\{\theta^\tau:\theta\in\operatorname{ElDiag}(\mathcal B)\}.
\]

\emph{Finite satisfiability.}  Let \(\Sigma_0\subseteq\Sigma\) be finite and
let \(\theta\) be the conjunction of the members of
\(\operatorname{ElDiag}(\mathcal B)\) whose translations lie in
\(\Sigma_0\); take \(\theta\) to be \(\forall x\,(x=x)\) if there are none.
Then \(\theta\in\operatorname{ElDiag}(\mathcal B)\), since a conjunction of
sentences true in \(\mathcal B_B\) is true there.  Write
\[
 \theta=\varphi(c_{a_{i_1}},\ldots,c_{a_{i_n}},c_{b_1},\ldots,c_{b_m}),
\]
where \(\varphi(x_{i_1},\ldots,x_{i_n},y_1,\ldots,y_m)\) is an
\(\Lang\)-formula, each occurring coordinate of \(\bar a\) is written as one
chosen \(a_{i_l}\), and \(b_1,\ldots,b_m\) are distinct elements of \(B\)
that are not coordinates of \(\bar a\).  Since
\(\mathcal B\entails\varphi(a_{i_1},\ldots,a_{i_n},b_1,\ldots,b_m)\),
\[
 \mathcal B\entails\exists y_1\cdots\exists y_m\,
   \varphi(a_{i_1},\ldots,a_{i_n},y_1,\ldots,y_m).
\]
This is an \(\Lang\)-formula of named elements, so by hypothesis
\(\mathcal C\mathbin{\upharpoonright}\Lang\) satisfies it at
\(c_{i_1},\ldots,c_{i_n}\), and by the reduct fact so does \(\mathcal C\).
Choose \(w_1,\ldots,w_m\in C\) with
\(\mathcal C\entails\varphi(c_{i_1},\ldots,c_{i_n},w_1,\ldots,w_m)\).
Expand \(\mathcal C_C\) to \(\Lang^*\) by interpreting \(d_{b_l}\) as
\(w_l\) for \(l\le m\) and every other \(d_b\) as an arbitrary element of
\(C\).  This expansion satisfies
\(\operatorname{ElDiag}(\mathcal C)\), whose sentences do not contain any
\(d_b\).  It satisfies \(\theta^\tau\), which is
\(\varphi(e_{c_{i_1}},\ldots,e_{c_{i_n}},d_{b_1},\ldots,d_{b_m})\), because
these constants denote \(c_{i_1},\ldots,c_{i_n},w_1,\ldots,w_m\).  Every
member of \(\Sigma_0\) coming from \(\mathcal B\) is a conjunct of
\(\theta^\tau\), so \(\Sigma_0\) has a model.

\emph{The maps.}  By \cref{thm:fol-compactness}, with the Axiom of Choice
when \(\Lang^*\) is uncountable, \(\Sigma\) has a model \(\mathcal D^*\).
Let \(\mathcal D_1=\mathcal D^*\mathbin{\upharpoonright}\Lang^+\).  The
\(\Lang^+(C)\)-reduct of \(\mathcal D^*\) satisfies
\(\operatorname{ElDiag}(\mathcal C)\), so by the second clause of
\cref{thm:diagram-tests}, in the signature \(\Lang^+\), the map
\(k(c)=e_c^{\mathcal D^*}\) is an elementary embedding
\(\mathcal C\to\mathcal D_1\).  Let \(\mathcal D_1'\) be the
\(\Lang(B)\)-expansion of \(\mathcal D_1\mathbin{\upharpoonright}\Lang\)
interpreting each \(c_b\) as \(\tau(b)^{\mathcal D^*}\).  For
\(\theta\in\operatorname{ElDiag}(\mathcal B)\), the sentence \(\theta\) in
\(\mathcal D_1'\) and the sentence \(\theta^\tau\) in \(\mathcal D^*\) are
evaluated with the same domain and the same interpretations of the symbols
they contain, so \(\mathcal D_1'\entails\theta\).  By the second clause of
\cref{thm:diagram-tests}, in the signature \(\Lang\), the map
\(g_1(b)=\tau(b)^{\mathcal D^*}\) is an elementary embedding
\(\mathcal B\to\mathcal D_1\mathbin{\upharpoonright}\Lang\), and
\(g_1(a_i)=e_{c_i}^{\mathcal D^*}=k(c_i)\).

\emph{Renaming.}  Apply the renaming paragraph to \(k\), obtaining
\(\rho\) and \(\mathcal D=\rho[\mathcal D_1]\) with
\(\mathcal C\preccurlyeq\mathcal D\).  Since \(\rho\) is an isomorphism of
\(\Lang^+\)-structures, it is also an isomorphism of the
\(\Lang\)-reducts, and \(g=\rho\circ g_1\) is an \(\Lang\)-elementary
embedding \(\mathcal B\to\mathcal D\mathbin{\upharpoonright}\Lang\) with
\(g(a_i)=\rho(k(c_i))=c_i\).
\end{proof}
\end{annexitem}

\begin{annexitem}{annex:p03-w03-p03-006-elem}{thm:p03-w03-p03-006-elementary-amalgamation}{Elementary amalgamation}
\begin{proof}
Clause~1 is \cref{thm:p03-w03-p03-006-expansion-amalgamation} with
\(\Lang^+=\Lang\), where \(\mathcal C\mathbin{\upharpoonright}\Lang\) is
\(\mathcal C\) itself.

For clause~2, index named tuples by the set \(A\): put
\(\bar a=(f(x))_{x\in A}\) in \(B\) and \(\bar c=(f'(x))_{x\in A}\) in
\(C\).  For every formula \(\varphi(x_1,\ldots,x_n)\) and
\(x_1,\ldots,x_n\in A\),
\begin{align*}
 \mathcal B\entails\varphi(f(x_1),\ldots,f(x_n))
 &\iff \mathcal A\entails\varphi(x_1,\ldots,x_n)\\
 &\iff \mathcal C\entails\varphi(f'(x_1),\ldots,f'(x_n)),
\end{align*}
because \(f\) and \(f'\) are elementary.  Hence
\((\mathcal B,\bar a)\equiv(\mathcal C,\bar c)\).  Clause~1 gives
\(\mathcal D\succcurlyeq\mathcal C\) and an elementary
\(g:\mathcal B\to\mathcal D\) with \(g(f(x))=f'(x)\) for all \(x\in A\).
Let \(h=g\) and let \(k\) be the inclusion \(C\subseteq D\), which is
elementary because \(\mathcal C\preccurlyeq\mathcal D\).  Then
\(h\circ f=k\circ f'\).
\end{proof}
\end{annexitem}

\begin{annexitem}{annex:p03-w03-p03-006-disjoint}{thm:p03-w03-p03-006-disjoint-amalgamation}{Disjoint amalgamation over an elementary substructure}
\begin{proof}
\emph{Naming.}  Use the constants \(e_c=c_c\) of \(\Lang(C)\) for the
elements of \(C\).  For \(b\in A\) put \(\tau(b)=e_b\); for
\(b\in B\setminus A\) let \(\tau(b)=d_b\), new and pairwise distinct
constants.  As \(B\cap C=A\), every element of \(A\) is named once, by
the same constant on both sides.  Define \(\theta^\tau\) as in the proof of
\cref{thm:p03-w03-p03-006-expansion-amalgamation} and put
\[
 \Sigma=\operatorname{ElDiag}(\mathcal C)
 \cup\{\theta^\tau:\theta\in\operatorname{ElDiag}(\mathcal B)\}
 \cup\{d_b\ne e_c:b\in B\setminus A,\ c\in C\setminus A\}.
\]

\emph{Finite satisfiability.}  Let \(\Sigma_0\subseteq\Sigma\) be finite.
Let \(b_1,\ldots,b_m\) be distinct elements of \(B\setminus A\) such that
every \(d_b\) occurring in \(\Sigma_0\) is among
\(d_{b_1},\ldots,d_{b_m}\), and let \(f_1,\ldots,f_k\) list the elements
\(c\in C\setminus A\) for which some inequality \(d_b\ne e_c\) lies in
\(\Sigma_0\).  As before, the members of \(\Sigma_0\) coming from
\(\mathcal B\) are conjuncts of a single \(\theta^\tau\) with
\[
 \theta=\varphi(c_{\bar a},c_{b_1},\ldots,c_{b_m})\in\operatorname{ElDiag}(\mathcal B),
\]
where \(\bar a\) is a finite tuple from \(A\) and
\(\varphi(\bar x,y_1,\ldots,y_m)\) is an \(\Lang\)-formula; unused
\(y_l\) are allowed.  We claim that
\[
 \mathcal C\entails\exists y_1\cdots\exists y_m\Bigl(\varphi(\bar a,\bar y)
   \land\bigwedge_{l\le m}\bigwedge_{j\le k}y_l\ne f_j\Bigr).
\]
Suppose not.  Let \(\psi(\bar x,z_1,\ldots,z_k)\) be the \(\Lang\)-formula
\[
 \forall y_1\cdots\forall y_m\Bigl(\varphi(\bar x,\bar y)
   \rightarrow\bigvee_{l\le m}\bigvee_{j\le k}y_l=z_j\Bigr),
\]
where an empty disjunction is the false sentence \(\neg\,\forall x\,(x=x)\).
Then \(\mathcal C\entails\psi(\bar a,f_1,\ldots,f_k)\), hence
\(\mathcal C\entails\exists\bar z\,\psi(\bar a,\bar z)\).  The parameters
\(\bar a\) lie in \(A\) and \(\mathcal A\preccurlyeq\mathcal C\), so
\(\mathcal A\entails\exists\bar z\,\psi(\bar a,\bar z)\); choose
\(a'_1,\ldots,a'_k\in A\) with
\(\mathcal A\entails\psi(\bar a,a'_1,\ldots,a'_k)\).  Since
\(\mathcal A\preccurlyeq\mathcal B\), also
\(\mathcal B\entails\psi(\bar a,a'_1,\ldots,a'_k)\).  But
\(\mathcal B\entails\varphi(\bar a,b_1,\ldots,b_m)\), so the disjunction
holds at \(b_1,\ldots,b_m\): it is nonempty and \(b_l=a'_j\in A\) for some
\(l\) and \(j\).  This contradicts \(b_l\in B\setminus A\), and proves the
claim.

Choose witnesses \(w_1,\ldots,w_m\in C\) for the claim.  Expand
\(\mathcal C_C\) by interpreting \(d_{b_l}\) as \(w_l\) and every other
\(d_b\) arbitrarily.  The expansion satisfies
\(\operatorname{ElDiag}(\mathcal C)\) and \(\theta^\tau\), which is
\(\varphi(e_{\bar a},d_{b_1},\ldots,d_{b_m})\).  It satisfies each
inequality \(d_{b_l}\ne e_{f_j}\) in \(\Sigma_0\), since \(w_l\ne f_j\).
So \(\Sigma_0\) has a model.

\emph{The maps.}  Compactness, \cref{thm:fol-compactness}, gives a model
\(\mathcal D^*\) of \(\Sigma\).  As in the proof of
\cref{thm:p03-w03-p03-006-expansion-amalgamation}, with
\(\Lang^+=\Lang\), the maps \(k(c)=e_c^{\mathcal D^*}\) and
\(g_1(b)=\tau(b)^{\mathcal D^*}\) are elementary embeddings of
\(\mathcal C\) and \(\mathcal B\) into
\(\mathcal D_1=\mathcal D^*\mathbin{\upharpoonright}\Lang\), and
\(g_1(a)=k(a)\) for \(a\in A\).

We show \(g_1(B)\cap k(C)=k(A)\).  The inclusion \(\supseteq\) holds
because \(g_1\) and \(k\) agree on \(A\).  Let \(g_1(b)=k(c)\).  If
\(b\in A\), then \(k(b)=k(c)\) and \(b=c\) by injectivity, so the common
value lies in \(k(A)\).  Suppose \(b\in B\setminus A\).  If
\(c\in C\setminus A\), the sentence \(d_b\ne e_c\) of \(\Sigma\) gives
\(g_1(b)\ne k(c)\), which is false.  So \(c\in A\), and
\(g_1(b)=k(c)=g_1(c)\) gives \(b=c\in A\) by injectivity of \(g_1\), again
false.  Hence only \(b\in A\) occurs.

Finally apply the renaming paragraph to \(k\), obtaining
\(\mathcal D=\rho[\mathcal D_1]\succcurlyeq\mathcal C\), and let
\(g=\rho\circ g_1\).  Then \(g\) is elementary, \(g(a)=\rho(k(a))=a\) for
\(a\in A\), and
\(g(B)\cap C=\rho(g_1(B)\cap k(C))=\rho(k(A))=A\), because \(\rho\) is a
bijection mapping \(k(C)\) onto \(C\).
\end{proof}
\end{annexitem}

\begin{annexitem}{annex:p03-w03-p03-006-family}{thm:p03-w03-p03-006-family-amalgamation}{Amalgamating a family}
\begin{proof}
\emph{Naming.}  Take constants \(e_x\) for \(x\in X\), and for each
\(j\in J\) and \(b\in B_j\setminus X\) a new constant \(d_{j,b}\), all
pairwise distinct.  Let \(\tau_j(x)=e_x\) for \(x\in X\) and
\(\tau_j(b)=d_{j,b}\) otherwise, and let
\(\Sigma_j=\{\theta^{\tau_j}:\theta\in\operatorname{ElDiag}(\mathcal B_j)\}\),
with \(\theta^{\tau_j}\) defined as before.  Let
\(\Sigma=\bigcup_{j\in J}\Sigma_j\).

\emph{Finite subfamilies.}  Every finite subset of \(\Sigma\) lies in
\(\bigcup_{j\in J_0}\Sigma_j\) for some finite nonempty \(J_0\subseteq J\),
using that \(J\) is nonempty.  We show by induction on \(\abs{J_0}\) that
each such union has a model.  If \(J_0=\{j\}\), interpret \(e_x\) as \(x\)
and \(d_{j,b}\) as \(b\) in \(\mathcal B_j\); this expansion satisfies
\(\Sigma_j\).  Now let \(J_0=J_1\cup\{j\}\) with \(j\notin J_1\) and
\(J_1\) nonempty, and let \(\mathcal D^*\) be a model of
\(\bigcup_{i\in J_1}\Sigma_i\).  Put
\(\mathcal D_0=\mathcal D^*\mathbin{\upharpoonright}\Lang\) and
\(\bar u=(e_x^{\mathcal D^*})_{x\in X}\).  Fix \(i\in J_1\).  As in the
proof of \cref{thm:p03-w03-p03-006-expansion-amalgamation}, the map
\(b\mapsto\tau_i(b)^{\mathcal D^*}\) is an elementary embedding
\(\mathcal B_i\to\mathcal D_0\) sending each \(x\in X\) to
\(e_x^{\mathcal D^*}\).  Therefore
\[
 (\mathcal B_j,\bar x)\equiv(\mathcal B_i,\bar x)\equiv(\mathcal D_0,\bar u).
\]
Clause~1 of \cref{thm:p03-w03-p03-006-elementary-amalgamation} gives
\(\mathcal D_1\succcurlyeq\mathcal D_0\) and an elementary embedding
\(g:\mathcal B_j\to\mathcal D_1\) with \(g(x)=e_x^{\mathcal D^*}\) for
\(x\in X\).  Expand \(\mathcal D_1\) by interpreting every constant
occurring in \(\bigcup_{i\in J_1}\Sigma_i\) as in \(\mathcal D^*\), each
\(d_{j,b}\) as \(g(b)\), and any remaining constant arbitrarily.  A
sentence \(\theta^{\tau_i}\) with \(i\in J_1\) says that an
\(\Lang\)-formula holds of certain elements of \(D_0\); it holds in
\(\mathcal D_0\), hence in \(\mathcal D_1\) because
\(\mathcal D_0\preccurlyeq\mathcal D_1\).  A sentence \(\theta^{\tau_j}\)
says that an \(\Lang\)-formula true in \(\mathcal B_j\) at certain
elements holds at their images under \(g\), since \(e_x\) denotes
\(g(x)\) and \(d_{j,b}\) denotes \(g(b)\); it holds because \(g\) is
elementary.  So \(\bigcup_{i\in J_0}\Sigma_i\) has a model.

\emph{Conclusion.}  By \cref{thm:fol-compactness}, \(\Sigma\) has a model
\(\mathcal D^*\).  Let \(\mathcal D=\mathcal D^*\mathbin{\upharpoonright}\Lang\)
and \(g_j(b)=\tau_j(b)^{\mathcal D^*}\).  As before, each \(g_j\) is an
elementary embedding \(\mathcal B_j\to\mathcal D\) by the second clause of
\cref{thm:diagram-tests}, and \(g_j(x)=e_x^{\mathcal D^*}\) for every
\(j\) and every \(x\in X\).
\end{proof}
\end{annexitem}

\begin{annexitem}{annex:p03-w03-p03-006-exist}{thm:p03-w03-p03-006-existential-amalgamation}{Existential amalgamation}
\begin{proof}
\emph{From 2 to 1.}  Let \(\varphi=\exists\bar y\,\chi\) with \(\chi\)
quantifier-free, and suppose
\(\mathcal C\entails\exists\bar y\,\chi(\bar c',\bar y)\), where
\(\bar c'\) lists the relevant coordinates of \(\bar c\) and \(\bar a'\)
the corresponding coordinates of \(\bar a\).  Choose witnesses \(\bar w\)
in \(C\).  The embedding \(g\) preserves quantifier-free formulas, so
\(\mathcal D\entails\chi(g\bar c',g\bar w)\), that is,
\(\mathcal D\entails\chi(\bar a',g\bar w)\).  Hence
\(\mathcal D\entails\exists\bar y\,\chi(\bar a',\bar y)\), and
\(\mathcal B\preccurlyeq\mathcal D\) gives
\(\mathcal B\entails\exists\bar y\,\chi(\bar a',\bar y)\).

\emph{From 1 to 2: naming.}  Use the constants \(e_b=c_b\) of
\(\Lang(B)\).  For \(c\in C\), let \(\tau(c)=e_{a_i}\) if \(c=c_i\); this
is independent of \(i\), since if \(c_i=c_j\), the quantifier-free, hence
existential, formula \(x_i=x_j\) holds of \(\bar c\) and so of \(\bar a\).
For other \(c\), let \(\tau(c)=d_c\), new and pairwise distinct.  Put
\[
 \Sigma=\operatorname{ElDiag}(\mathcal B)
   \cup\{\theta^\tau:\theta\in\operatorname{Diag}(\mathcal C)\}.
\]

\emph{Finite satisfiability.}  Let \(\Sigma_0\subseteq\Sigma\) be finite,
and let \(\theta\) be the conjunction of the members of
\(\operatorname{Diag}(\mathcal C)\) whose translations occur in it,
or the atomic sentence \(c_c=c_c\) for some \(c\in C\) if there are none.
It is a quantifier-free sentence true in \(\mathcal C_C\).  Write
\(\theta=\chi(c_{c_{i_1}},\ldots,c_{c_{i_n}},c_{w_1},\ldots,c_{w_m})\), with
\(\chi(x_{i_1},\ldots,x_{i_n},y_1,\ldots,y_m)\) a quantifier-free
\(\Lang\)-formula, the \(w_l\) distinct and not coordinates of \(\bar c\).
Then \(\mathcal C\entails\exists\bar y\,\chi(c_{i_1},\ldots,c_{i_n},\bar y)\),
an existential formula of named elements, and the hypothesis gives
\(\mathcal B\entails\exists\bar y\,\chi(a_{i_1},\ldots,a_{i_n},\bar y)\)
with witnesses \(v_1,\ldots,v_m\) in \(B\).  Expand \(\mathcal B_B\) by
interpreting \(d_{w_l}\) as \(v_l\), and the remaining \(d_c\)
arbitrarily.  The expansion satisfies \(\operatorname{ElDiag}(\mathcal B)\)
and \(\theta^\tau=\chi(e_{a_{i_1}},\ldots,e_{a_{i_n}},d_{w_1},\ldots,d_{w_m})\),
so \(\Sigma_0\) has a model.

\emph{The maps.}  Compactness gives a model \(\mathcal D^*\) of
\(\Sigma\); let \(\mathcal D_1=\mathcal D^*\mathbin{\upharpoonright}\Lang\).
By the second clause of \cref{thm:diagram-tests}, \(k(b)=e_b^{\mathcal D^*}\)
is an elementary embedding \(\mathcal B\to\mathcal D_1\).  The
\(\Lang(C)\)-expansion of \(\mathcal D_1\) interpreting \(c_c\) as
\(\tau(c)^{\mathcal D^*}\) satisfies \(\operatorname{Diag}(\mathcal C)\),
by the same comparison of symbols used in the proof of
\cref{thm:p03-w03-p03-006-expansion-amalgamation}; by the first clause of
\cref{thm:diagram-tests}, \(g_1(c)=\tau(c)^{\mathcal D^*}\) is an
embedding \(\mathcal C\to\mathcal D_1\), and
\(g_1(c_i)=e_{a_i}^{\mathcal D^*}=k(a_i)\).  Apply the renaming paragraph
to \(k\), obtaining \(\mathcal D=\rho[\mathcal D_1]\succcurlyeq\mathcal B\),
and let \(g=\rho\circ g_1\).  Then \(g\) is an embedding and
\(g(c_i)=\rho(k(a_i))=a_i\).
\end{proof}
\end{annexitem}

\begin{annexitem}{annex:p03-w03-p03-006-univ}{thm:p03-w03-p03-006-universal-part}{Models of the universal part}
\begin{proof}
Suppose \(g:\mathcal B\to\mathcal M\) is an embedding with
\(\mathcal M\entails T\), and let
\(\forall\bar y\,\chi(\bar y)\in T_\forall\) with \(\chi\)
quantifier-free.  For every tuple \(\bar b\) from \(B\),
\(\mathcal M\entails\chi(g\bar b)\), and since embeddings reflect
quantifier-free formulas, \(\mathcal B\entails\chi(\bar b)\).  Hence
\(\mathcal B\entails\forall\bar y\,\chi(\bar y)\).

Conversely, suppose \(\mathcal B\entails T_\forall\) and let
\(\Sigma=T\cup\operatorname{Diag}(\mathcal B)\) in \(\Lang(B)\).  Let
\(\Sigma_0\subseteq\Sigma\) be finite, and let \(\chi(c_{b_1},\ldots,c_{b_m})\)
be the conjunction of its members from \(\operatorname{Diag}(\mathcal B)\),
or the atomic sentence \(c_b=c_b\) for some \(b\in B\) if there are none,
with \(\chi(y_1,\ldots,y_m)\)
quantifier-free and \(b_1,\ldots,b_m\) distinct.  Suppose
\(T\cup\{\chi(c_{b_1},\ldots,c_{b_m})\}\) had no model.  Let
\(\mathcal M\entails T\) and \(m_1,\ldots,m_m\in M\).  The expansion of
\(\mathcal M\) interpreting \(c_{b_l}\) as \(m_l\) still satisfies \(T\),
whose sentences do not contain these constants, so it satisfies
\(\neg\chi(c_{b_1},\ldots,c_{b_m})\); that is,
\(\mathcal M\entails\neg\chi(m_1,\ldots,m_m)\).  Hence
\(T\entails\forall\bar y\,\neg\chi(\bar y)\).  This universal sentence lies
in \(T_\forall\), so \(\mathcal B\) satisfies it, contradicting
\(\mathcal B\entails\chi(b_1,\ldots,b_m)\).  Therefore
\(T\cup\{\chi(c_{\bar b})\}\) has a model, and so does \(\Sigma_0\).

By \cref{thm:fol-compactness}, \(\Sigma\) has a model \(\mathcal M^*\).
Its reduct \(\mathcal M=\mathcal M^*\mathbin{\upharpoonright}\Lang\)
satisfies \(T\) by the reduct fact, and by the first clause of
\cref{thm:diagram-tests} the map \(b\mapsto c_b^{\mathcal M^*}\) is an
embedding \(\mathcal B\to\mathcal M\).
\end{proof}
\end{annexitem}

\begin{annexitem}{annex:p03-w03-p03-006-reduct}{thm:p03-w03-p03-006-reduct-models}{Models of the restricted theory}
\begin{proof}
Suppose \(\mathcal M\entails T\) and
\(g:\mathcal B\to\mathcal M\mathbin{\upharpoonright}\Lang\) is
\(\Lang\)-elementary.  For \(\sigma\in T^{\Lang}\),
\(\mathcal M\entails\sigma\), so
\(\mathcal M\mathbin{\upharpoonright}\Lang\entails\sigma\) by the reduct
fact.  An elementary embedding preserves and reflects sentences, which are
formulas without free variables, so \(\mathcal B\entails\sigma\).

For the converse, we first prove the moreover clause.  Let
\(\mathcal M_0\entails T\) with
\(\mathcal B\equiv\mathcal M_0\mathbin{\upharpoonright}\Lang\).  Apply
\cref{thm:p03-w03-p03-006-expansion-amalgamation} with
\(\mathcal C=\mathcal M_0\) and \(I=\varnothing\); its hypothesis is then
exactly \(\mathcal B\equiv\mathcal M_0\mathbin{\upharpoonright}\Lang\).  It
yields an \(\Lang^+\)-structure \(\mathcal M\succcurlyeq\mathcal M_0\) and
an \(\Lang\)-elementary embedding
\(\mathcal B\to\mathcal M\mathbin{\upharpoonright}\Lang\).  An elementary
extension satisfies the same sentences, so \(\mathcal M\entails T\).

Now suppose only \(\mathcal B\entails T^{\Lang}\), and let
\(\operatorname{Th}(\mathcal B)\) be the set of \(\Lang\)-sentences true in
\(\mathcal B\).  We show that \(T\cup\operatorname{Th}(\mathcal B)\) has a
model.  Otherwise, by \cref{thm:fol-compactness}, some finite
\(T_0\subseteq T\) and \(\sigma_1,\ldots,\sigma_n\in\operatorname{Th}(\mathcal B)\)
have no common model.  Let \(\sigma\) be \(\sigma_1\land\cdots\land\sigma_n\),
or \(\forall x\,(x=x)\) when \(n=0\).  Every model of \(T\) then satisfies
\(\neg\sigma\), so \(\neg\sigma\in T^{\Lang}\) and
\(\mathcal B\entails\neg\sigma\); but \(\mathcal B\entails\sigma\).  Let
\(\mathcal M_0\entails T\cup\operatorname{Th}(\mathcal B)\).  For an
\(\Lang\)-sentence \(\sigma\), if \(\mathcal B\entails\sigma\) then
\(\sigma\in\operatorname{Th}(\mathcal B)\) and
\(\mathcal M_0\mathbin{\upharpoonright}\Lang\entails\sigma\); if
\(\mathcal B\not\models\sigma\) then
\(\neg\sigma\in\operatorname{Th}(\mathcal B)\) and
\(\mathcal M_0\mathbin{\upharpoonright}\Lang\not\models\sigma\).  So
\(\mathcal B\equiv\mathcal M_0\mathbin{\upharpoonright}\Lang\), and the
moreover clause gives the required \(\mathcal M\).
\end{proof}
\end{annexitem}

\begin{annexitem}{annex:p03-w03-p03-006-finite-avoidance}{thm:p03-w03-p03-006-finite-avoidance}{Finite covers of a definable family of tuples}
\begin{proof}
We first prove a finite-set fact. A set \(S\) \emph{meets} a family
\(H\) of nonempty finite sets if \(S\cap D\neq\varnothing\) for
every \(D\in H\). If no set of size at most \(k\) meets \(H\), then
some finite subfamily already has this property. Induct on the
nonnegative integer \(k\). For \(k=0\), choose any member of \(H\).
For \(k>0\), choose \(D\in H\). For each \(d\in D\), the family
\(H_d=\{E\in H:d\notin E\}\) cannot be met by at most \(k-1\)
points: adjoining \(d\) would then meet all of \(H\) with at most
\(k\). By induction choose a finite subfamily \(H'_d\subseteq H_d\)
that cannot be met by \(k-1\) points. The finite family
\(\{D\}\cup\bigcup_{d\in D}H'_d\) cannot be met by \(k\) points.
Indeed any putative meeting set \(S\) meets \(D\) at some \(d\), and
\(S\setminus\{d\}\) still meets every member of \(H'_d\), since
none contains \(d\). This contradicts its choice.

Now let \(H\) consist of the coordinate sets
\(\{c_1,\ldots,c_m\}\) of all tuples satisfying
\(\varphi(\bar y,\bar a)\) in \(\mathcal M\). These are nonempty
finite sets, and \(H\) is nonempty by the given tuple. Suppose some
finite set meets \(H\), and let \(r\geq1\) be the least cardinality
of such a set. No set of size at most \(r-1\) meets \(H\). The fact
just proved supplies a finite \(H_0\subseteq H\) with the same
property. Put \(U=\bigcup H_0\); this is finite. Every size-\(r\)
set meeting \(H\) is contained in \(U\), since deleting any member
outside \(U\) would leave at most \(r-1\) points meeting \(H_0\).

The elements belonging to a size-\(r\) set meeting \(H\) form a set
definable over \(X\). A formula for membership of \(z_1\) is
\[
 \exists z_2\cdots\exists z_r\left(
   \bigwedge_{i<j}z_i\neq z_j\ \land
   \forall y_1\cdots\forall y_m
   \left(\varphi(\bar y,\bar a)\to
                  \bigvee_{i\leq m,\,j\leq r}y_i=z_j\right)
                                      \right).
\]
For \(r=1\), omit the existential block and read the empty conjunction
as truth. The defined set is contained in finite \(U\), and contains
every point of every minimum meeting set. Hence all these points belong
to \(\operatorname{acl}_{\mathcal M}(X)\). Fix one minimum meeting
set \(S\). It meets the coordinate set of the given tuple \(\bar b\),
so some \(b_i\) belongs to \(S\) and is algebraic over \(X\), a
contradiction. Thus no finite set meets every satisfying tuple, which is
exactly the desired avoidance conclusion.
\end{proof}
\end{annexitem}

\begin{annexitem}{annex:p03-w03-p03-006-acl-amalgamation}{thm:p03-w03-p03-006-acl-amalgamation}{Disjointness over an algebraically closed common set}
\begin{proof}
Use one constant \(e_x\) for each \(x\in X\), additional constants
\(e_c\) for \(c\in C\setminus X\), and fresh constants \(d_b\) for
\(b\in B\setminus X\). Translate the elementary diagram of
\(\mathcal B\) with \(x\in X\) named by \(e_x\) and each remaining
\(b\) by \(d_b\). Let \(\Sigma\) be that translated diagram, the
elementary diagram of \(\mathcal C\), and every inequality
\(d_b\neq e_c\) with \(b\in B\setminus X\),
\(c\in C\setminus X\).

Fix a finite subset of \(\Sigma\). Conjoin its sentences from the
\(\mathcal B\)-diagram and replace the finitely many distinct new
constants by variables. This gives \(\varphi(\bar y,\bar a)\), with
parameters \(\bar a\) from \(X\), true in \(\mathcal B\) at a tuple
\(\bar b\) from \(B\setminus X\). Include unused variables for
constants appearing only in inequalities. If this tuple is nonempty,
every coordinate lies outside \(\operatorname{acl}_{\mathcal B}(X)\).
For each finite \(k\), the avoidance lemma gives
\[
 \mathcal B\entails
 \forall z_1\cdots\forall z_k\exists\bar y
 \left(\varphi(\bar y,\bar a)\land
                  \bigwedge_{i,j}y_i\neq z_j\right).
\]
This is a formula over the common parameter set, so it holds in
\(\mathcal C\) too. Take the \(z_j\) to list every element of
\(C\setminus X\) forbidden by an inequality in the chosen finite
subset. The resulting tuple in \(C\) interprets the \(d_b\)'s,
satisfying both the translated \(\mathcal B\)-sentences and all those
inequalities. Interpreting each \(e_c\) as \(c\) keeps the entire
\(\mathcal C\)-diagram true. If \(\bar b\) is empty, no inequalities
of that form occur, and common-parameter agreement transfers
\(\varphi(\bar a)\) directly. This proves finite satisfiability.

By compactness, \(\Sigma\) has a model. Its naming maps are elementary
embeddings of \(\mathcal B\) and \(\mathcal C\), agreeing on \(X\).
The added inequalities forbid intersections between the images of
\(B\setminus X\) and \(C\setminus X\); injectivity also forbids an
image from either of these sets equalling the image of an element of
\(X\). Their images therefore intersect exactly in the image of
\(X\). The renaming construction at the start of this annex makes the
\(\mathcal C\)-embedding the inclusion, giving the stated
\(\mathcal D\) and \(g\).

For the last assertion, let \(\theta(z,\bar a)\) over \(X\) have
exactly \(r\) solutions in \(\mathcal C\). The same first-order
counting assertion holds in \(\mathcal B\). All its solutions there
belong to \(X\), by algebraic closure. Name them \(a_1,\ldots,a_r\).
The formula
\(\forall z(\theta(z,\bar a)\leftrightarrow
\bigvee_{i\leq r}z=a_i)\) is over \(X\), hence transfers to
\(\mathcal C\). Its solutions in \(\mathcal C\) are exactly those
same elements of \(X\). The case \(r=0\) is vacuous. Thus
\(\operatorname{acl}_{\mathcal C}(X)=X\).
\end{proof}
\end{annexitem}
